The signed operator is designed based on the assumption that the edge semantic can be represented under a signed, continuous value rather than a binary decision.
Given the predicted heterophily score $\hat{q}_{uv}$ in Phase 1, we start by constructing the signed adjacency matrix $\hat{S} = [\hat{s}_{uv}]$ 
\begin{equation}
    \hat{s}_{uv} = 1 - 2 \hat{q}_{uv},
    \label{eq:suv}
\end{equation}
with $\hat{s}_{uv} \in [-1,1]$.
Note that if we know the true value $s_{uv}$ we can easily reconstruct the underlying labels from the training nodes.
However, due to the lack of node labels in semi-supervised settings, we only have a noisy estimate of $s_{uv}$ through $\hat{q}_{uv}$.

To incorporate the uncertainty of heterophilic edge semantics, we introduce a
learned parameter $\eta\in[0,1]$ as the signed weight applied to the self-looped
adjacency matrix $\tilde{\bm A}$,
\begin{equation}
    \bm{W}^{(\eta)} \;=\; \tilde{\bm{A}}\odot\bigl[(1-\eta)\bm{1}\bm{1}^{\!\top}
    + \eta\,\hat{\bm{S}}\bigr],
    \label{eq:weta}
\end{equation}
where $\hat{\bm S}=\bm{1}\bm{1}^{\!\top}-2\hat{\bm Q}$ is the soft sign matrix
induced by the frozen probabilities $\hat{\bm Q}=[\hat q_{uv}]$,
with $\hat q_{uv}=\hat q_{vu}\in[0,1]$ and $\hat q_{uu}\equiv 0$.
Expanding Eq.~\eqref{eq:weta} entrywise gives the scalar form
\begin{equation}
W^{(\eta)}_{uv} \;=\; \tilde A_{uv}\,\bigl(1-2\eta\,\hat q_{uv}\bigr).
\label{eq:Wentry}
\end{equation}
The role of calibration strength $\eta$ is a global gain on the signed correction.
When edge predictions are unreliable, $\eta$ can shrink to $0$ and the operator reduces to the unsigned Laplacian of BWGNN.
ESO-GAD therefore extends BWGNN along the operator axis, providing the same filter bank with a spectrum that carries calibrated edge semantics.
Moreover, the same scalar controls where signs flip.
Since $W^{(\eta)}_{uv}<0$ requires $\hat q_{uv}>1/(2\eta)$, and
since $\hat q_{uv}\le 1$, no edge can change sign unless $\eta>\tfrac12$.
For $\eta\ge\tfrac12$, an edge is exactly pruned when $\hat q_{uv}=1/(2\eta)$.
Because Eq.~\eqref{eq:Wentry} is affine in
$\hat q_{uv}$, no choice of $\eta$ reproduces a thresholded pruning rule, and edge deletion survives only on the measure-zero level set $\{\hat q_{uv}=1/(2\eta)\}$.
This distinguishes ESO-GAD from hard-pruning approaches to heterophily \citep{GHRN, CARE-GNN}, which discard inter-class edges and therefore
receive no gradient signal from them once removed.
By contrast, training $\eta$ jointly with the filters lets the model invert, or restore the contribution of a predicted-heterophilic edge, and every edge retains a gradient path throughout optimization.
Besides the signed weight calibration in Eq.~\eqref{eq:weta}, reweighting approaches such as \citep{SEC-GFD} can be used to construct the operator.
We provide a comparison with pruning and reweighting approaches in the experimental section.

The signed weights $W^{(\eta)}_{uv}$ in Eq.~\eqref{eq:Wentry} break the usual degree normalization due to the possible negative values.
Therefore, we normalize the degree by using the absolute degree \citep{KunegisSigned}:
\begin{equation}
D^{(\eta)}_{uu}=\textstyle\sum_v \bigl|W^{(\eta)}_{uv}\bigr|,
\label{eq:degreelaplacian}
\end{equation}
and create the Laplacian accordingly
\begin{equation}
\Lb \;=\; \bm{I}-\Db{}^{-1/2}\,\Wb\,\Db{}^{-1/2}.
\label{eq:laplacian}
\end{equation}
Its quadratic form decomposes edge-wise as
\begin{equation}
\bm{h}^{\!\top}\Lb\bm{h}
=\tfrac12\sum_{u,v}\bigl|W^{(\eta)}_{uv}\bigr|
\Bigl(\tfrac{h_u}{\sqrt{D^{(\eta)}_{uu}}}-\sgn\!\bigl(W^{(\eta)}_{uv}\bigr)\tfrac{h_v}{\sqrt{D^{(\eta)}_{vv}}}\Bigr)^{\!2},
\label{eq:quadform}
\end{equation}
which smooths across positively signed edges and enforces contrast across negatively signed ones, with per-edge strength $|1-2\eta \hat{q}_{uv}|$.
Heterophilic edges from the calibration phase are retained with reversed semantics rather than discarded like other pruning-based approaches.

Following BWGNN~\citep{BWGNN}, the $k$-th channel applies the Beta-wavelet kernel to $\Lb$, with the spectrum mapped from $[0,2]$ to the Beta support $[0,1]$ as:
\begin{equation}
\Wk(\Lb)
=\frac{1}{2\,\mathrm B(k{+}1,\,P{-}k{+}1)}
\Bigl(\tfrac{\Lb}{2}\Bigr)^{\!k}
\Bigl(\bm I-\tfrac{\Lb}{2}\Bigr)^{\!P-k},
\label{eq:betakernel}
\end{equation}
for $k=0,\dots,P$, giving $P+1$ channels in total.
In Eq.~\eqref{eq:betakernel}, $\mathrm B(\cdot,\cdot)$ is the Beta function, $P$ is the wavelet order and
the factor $\tfrac12$ is the Jacobian of the rescaling.
Under this formulation, the channel $k=0$  and $k=P$ are the low-pass filter and the high-pass filter, respectively.
Intermediate channels cover band-pass responses whose peaks sweep monotonically across the calibrated spectrum.
It is required that the kernel $\Wk(\Lb)$ is valid for any $\Lb$ constructed from the learned $\hat{q}_{uv}$ and a valid $\eta$, which is established by Proposition~\ref{prop:valid}.
\begin{proposition}[Spectral validity]\label{prop:valid}
Let $\tilde{\bm A}=\bm A+\bm I$ and let the scores be symmetric with $\hat q_{uu}=0$. For every such $\hat q\in[0,1]^{|\mathcal E|}$ and every $\eta\ge0$, in particular for every $\eta\in[0,1]$, the following conditions hold
\[
\Db\succeq \bm I \succ 0
\qquad\text{and}\qquad
\lambda_i\in[0,2],
\]
where $\lambda_i$ is an eigenvalue of $\Lb$.
\end{proposition}
%As a result, every Beta-wavelet kernel $\Wk(\Lb)$ of Eq.~\eqref{eq:betakernel} is well defined
In implementation, the product form in Eq.~\eqref{eq:betakernel} is not computed directly.
Instead, each kernel is pre-expanded into monomials of $\Lb$ with fixed coefficients $\theta^{(k)}_j$,
\begin{equation}
\Wk(\Lb)=\sum_{j=0}^{P}\theta^{(k)}_j\,\Lb^{\,j}.
\label{eq:monomial}
\end{equation}
At the first layer, a MLP network with parameters $\psi$ is used to map the input feature space into the initial hidden representation $\bm H^{(0)}= \text{MLP}_\psi(\bm X)$.
Then, every channel
\begin{equation}
\bm H_k=\Wk(\Lb)\,\bm H^{(0)},\qquad k=0,\dots,P,
\label{eq:channels}
\end{equation}
is computed using the fixed coefficients in Eq.~\eqref{eq:monomial}.
%The total computation is evaluated by at most $D$ sparse matrix-vector products
%The powers $\Lb^{\,j}\bm H^{(0)}$ are shared across channels, so the whole bank costs $O(D\,|\mathcal E|\,d)$ per forward pass.

Because $\eta$ is trained jointly with the filters, Eq.~\eqref{eq:weta} and Eq.~\eqref{eq:laplacian} are re-evaluated at every forward pass, which makes the filter bank a moving operator, and only at convergence is there one fixed $\Lb$ per dataset.
A small update of $\eta$ would introduce stochasticity into the underlying operator.
In order to improve the stability of the training process, we design an anchored gate in the later weighted aggregation $\bm H_k$ to produce the output latent representation.
An additional KL loss is added to prevent the collapse issue when the signal is not informative at the start of the training phase.
Specifically, the channel outputs are combined by a learnable simplex gate $\bm{g}=\mathrm{softmax}(\bm{a})$:
\begin{equation}
\bm{Z}=\big\Vert_{k=0}^{P}\sqrt{g_k}\,\bm{H}_k,
\label{eq:gate}
\end{equation}
 where $\bm{a}$ is trainable.
The classification output head $\text{MLP}_\omega$ is then used to map to the target label as:
\begin{equation}
\hat{\bm{y}}= \text{MLP}_\omega(\bm{Z}).
\end{equation}
Given the set of available labels $\bm{y}^L$ and the predicted output $\hat{\bm{y}}^L$, the whole module is  trained with
\begin{equation}
\mathcal{L}
=\mathcal{L}_{\mathrm{BCE}}(\bm{y}^L,\hat{\bm{y}}^L)
+\alpha\, D_{\mathrm{KL}}\!\bigl(\bm{g}\,\Vert\,\bm{g}_0\bigr),
\label{eq:loss}
\end{equation}
where $\bm{g}_0$ is the uniform prior and $\alpha$ is the weighting coefficient.
The gate acts as an adaptive frequency response for readout procedure.
The learned $\bm{g}$ decides which signed-spectral bands carry anomaly evidence on a given dataset, replacing the standard Beta-wavelet's fixed concatenation.
By combining the gate with the KL loss, it stabilizes training under a moving operator.
The pseudocode of ESO-GAD is presented in Algorithm~\ref{alg:ecsbwgnn}.

The runtime complexity of ESO-GAD combines two phases.
In the first phase, with $n=|\mathcal V|$, $m=|\mathcal E|$ and edge-descriptor dimension $d_e$, the training cost for each fold is approximately $O(m\,d_e)$, which is linear in the number of edges.
In Phase~2, each forward pass rebuilds $\bm W^{(\eta)}$ and
$\bm D^{(\eta)}$ from the frozen $\hat{\bm Q}$ in $O(m)$ time and shares the powers $\bm L_\eta^{\,j}\bm H^{(0)}$ across all $P{+}1$ channels via Eq.~\eqref{eq:monomial}.
This results in a per-epoch cost of
$O\bigl(P\,m\,d_h+P^{2}n\,d_h\bigr)$ with the hidden width $d_h$ of both $\psi$ and $\omega$.
The cost matches the standard Beta-wavelet bank~\citep{BWGNN}.
The coefficients $\theta^{(k)}_j$ are constants that do not depend on $\eta$; what changes with $\eta$ is the operator, which is rebuilt at every step in $O(m)$ time, whereas the unsigned bank builds its operator once. This additional cost is dominated by the $O(P\,m\,d_h)$ propagation by a factor of $P\,d_h$.
